\section{Discussion}

\subsection{What the Contract Reveals About Substrates}

Building three adaptation layers made the contract a diagnostic instrument as
well as an interface. The distribution of adaptation effort is not uniform:
memory and interrupt operations demanded the most bridging and recovery on every
non-trivial substrate, while synchronization and lifecycle operations mapped
almost directly. This pattern is consistent across \asterinas{} and Linux
despite their architectural differences, which suggests that the contract's
structure---not any single OS's shape---identifies where OS designs genuinely
diverge. A substrate that makes physical memory and stage-2 control cheap to
obtain (like \arceos{}) is hypervisor-friendly almost for free; a substrate that
abstracts them away (like Linux) pays in recovery. The contract therefore
doubles as a specification of what it means for an OS to be a good hypervisor
substrate.

\subsection{Limitations of the Current Contract}

The contract is sufficient for our guest workloads but not yet complete. Device
passthrough depends on optional IOMMU capabilities that not all substrates
expose, so full passthrough is substrate-gated. Multi-core guest scheduling and
CPU affinity are mediated by the substrate's scheduler, which the contract
treats as policy; richer control over vCPU placement would require extending the
contract. Live migration, snapshot/restore, and nested virtualization are out of
scope. These are natural extensions rather than refutations: each would add
operations to the contract, but the core's substrate-independence would be
preserved because the additions are themselves substrate-neutral in
specification.

\subsection{Generality Beyond the Three Substrates}

The three substrates span unikernel, framekernel, and monolithic designs, which
covers much of the practical OS design space for our target domains. Microkernel
and capability-based substrates (e.g., seL4-based systems) remain untested. We
conjecture that such substrates are amenable to the contract because their
service-oriented structure already separates policy from mechanism, but
validating this is future work. The contract's eligibility criterion---a
substrate can host \axvisor{} iff it provides the required capabilities---makes
the question empirical rather than architectural: any OS that can allocate
physical memory, control stage-2 translation, create execution contexts, and
route interrupts is a candidate.

\subsection{Implications for OS and Hypervisor Co-Design}

The work suggests a shift in how virtualization support is added to new OSes.
Rather than treating the hypervisor as an OS-specific artifact built during
kernel development, a community could maintain a shared, portable hypervisor
core and treat each OS's adaptation layer as a bounded engineering task
comparable to writing a driver. This decouples the evolution of virtualization
logic from the evolution of any single OS and lets emerging OSes inherit mature
virtualization by contract compliance. The KVM-compatible ABI amplifies this
benefit on the guest side: the same guest ecosystem is inherited for free. The
net effect is to reduce the $O(\text{OSes} \times \text{guests})$ implementation
matrix toward $O(\text{OSes}) + O(\text{guests})$, mediated by one core and two
stable contracts.
